% Source files for every number in this section (2026-07-31):
%   labels           /mnt/scratch/lh/labels/final3_env_rc_*.npz
%   sigma sweep      /mnt/scratch/lh/labels/final3_env_sig{025,100}_rc_*.npz
%   heterogeneity    results/analysis/heterogeneity_claim1.txt
%   shadow/online    results/robocasa/switch_controls/<task>_shadow_16.json
%   fixed k          results/robocasa/ksweep/<task>_fixed_<k>.json
%   ACT              results/robocasa/act/<task>_act_fixed_<k>.json
%   oracle           results/robocasa/oracle/
%   predictors       results/predictor/rc_seq_v3/S_<task>_h<1|4>_s<0|1|2>/{metrics,calib}.json
%   online switching results/robocasa/seqpred_calib/, thresholds in
%                    results/predictor/rc_seq_v3/online_thresholds.json
%   fork             results/robocasa/fork/<task>_fork_16-1.json
%   proxies          results/robocasa/proxy/<task>_proxyshadow_16.json
%   diffusion policy results/stage2/ksweep/, results/stage2/ksweep_libero/, results/ksweep/

\section{Claim 1: One Task Contains Both Regimes}
\label{sec:claim1}

\subsection{Both regimes coexist inside single episodes}

Table~\ref{tab:labels} summarizes the RoboCasa labels. Unstable stamps are a
small minority everywhere, between 2 and 5 percent of all timesteps, but they
are not spread evenly across episodes. Between 78 and 96 percent of individual
episodes contain both regimes, and a typical episode crosses between them
several times. The same procedure on LIBERO, MimicGen and robomimic gives the
same picture on twenty more datasets, with unstable fractions from 4 to 51
percent, both regimes present in essentially every episode, and 3 to 43
crossings per episode. Error therefore does not compound at a uniform rate
along an episode, which is what makes one chunk length per task the wrong unit.

\begin{table}[t]
\centering
\small
\begin{tabular}{lrrrrr}
\toprule
Task & stamps & $\delta$ & unstable & episodes with both & crossings \\
\midrule
OpenCabinet & 34{,}133 & .085 & 2.3\% & 86\% & 4 \\
OpenDrawer & 24{,}293 & .076 & 3.1\% & 78\% & 2 \\
PickPlaceCounterToCabinet & 23{,}724 & .065 & 5.0\% & 96\% & 6 \\
TurnOnSinkFaucet & 20{,}198 & .058 & 3.9\% & 82\% & 2 \\
\bottomrule
\end{tabular}
\caption{Open-loop stability labels on RoboCasa, 200 demonstrations per task.
$\delta$ is the automatic per-task threshold and unstable is the fraction of
stamps with $\lambda > \delta$. The last two columns show the two regimes
coexisting inside single episodes: the fraction of episodes containing both,
and the median number of crossings between them per episode.}
\label{tab:labels}
\end{table}

\begin{table*}[t]
\centering
\scriptsize
\setlength{\tabcolsep}{4pt}
\begin{tabular}{lrrrrrr}
\toprule
& & & \multicolumn{2}{c}{regime share, all stamps} & \multicolumn{2}{c}{heterogeneity within episodes} \\
\cmidrule(lr){4-5} \cmidrule(lr){6-7}
Dataset & stamps & $\delta$ & deadband & unstable & within-ep var & crossings/100 \\
\midrule
\multicolumn{7}{l}{\textit{RoboCasa, GR00T N1.5}} \\
\quad OpenCabinet & 34,133 & 0.085 & 95\% & 2.3\% & 99.4\% & 2.3 \\
\quad OpenDrawer & 24,293 & 0.076 & 94\% & 3.1\% & 97.8\% & 2.4 \\
\quad PickPlaceToCabinet & 23,724 & 0.065 & 93\% & 5.0\% & 97.8\% & 5.2 \\
\quad TurnOnSinkFaucet & 20,198 & 0.058 & 94\% & 3.9\% & 94.5\% & 2.8 \\
\addlinespace[2pt]
\multicolumn{7}{l}{\textit{LIBERO-10, diffusion policy}} \\
\quad K3 stove + moka & 6,060 & 0.085 & 93\% & 4.4\% & 96.9\% & 4.1 \\
\quad K4 bowl in drawer & 5,632 & 0.024 & 68\% & 29.8\% & 98.0\% & 9.9 \\
\quad K6 mug in microwave & 7,031 & 0.015 & 47\% & 50.9\% & 96.6\% & 22.8 \\
\quad K8 both moka pots & 9,812 & 0.014 & 47\% & 50.1\% & 98.1\% & 21.5 \\
\quad L1 soup + cream & 6,147 & 0.025 & 73\% & 25.0\% & 97.6\% & 12.2 \\
\quad L2 soup + tomato & 6,761 & 0.014 & 48\% & 49.8\% & 97.0\% & 21.3 \\
\quad L2 cream + butter & 5,920 & 0.014 & 48\% & 50.3\% & 97.0\% & 25.1 \\
\quad L5 mug + chocolate & 5,868 & 0.018 & 54\% & 43.1\% & 98.5\% & 11.5 \\
\quad L6 mug + chocolate & 5,789 & 0.018 & 60\% & 37.7\% & 96.5\% & 14.6 \\
\quad S1 book on shelf & 4,146 & 0.011 & 47\% & 50.2\% & 91.9\% & 15.8 \\
\addlinespace[2pt]
\multicolumn{7}{l}{\textit{MimicGen, diffusion policy}} \\
\quad coffee preparation & 66,724 & 0.064 & 85\% & 12.2\% & 99.6\% & 9.2 \\
\quad nut assembly & 33,251 & 0.056 & 83\% & 14.8\% & 98.8\% & 12.8 \\
\quad square & 12,874 & 0.036 & 61\% & 36.3\% & 95.5\% & 17.5 \\
\quad stack & 8,445 & 0.022 & 55\% & 42.7\% & 95.7\% & 23.2 \\
\addlinespace[2pt]
\multicolumn{7}{l}{\textit{robomimic, diffusion policy}} \\
\quad can & 9,254 & 0.050 & 83\% & 14.2\% & 95.9\% & 12.8 \\
\quad lift & 2,485 & 0.018 & 55\% & 42.1\% & 95.5\% & 23.2 \\
\quad square & 12,723 & 0.037 & 64\% & 33.4\% & 94.8\% & 15.2 \\
\quad tool\_hang & 45,633 & 0.035 & 57\% & 40.3\% & 98.2\% & 19.1 \\
\addlinespace[2pt]
\multicolumn{7}{l}{\textit{policy rollouts (not demonstrations)}} \\
\quad can (rollouts) & 8,360 & 0.038 & 66\% & 31.6\% & 89.7\% & 17.5 \\
\quad lift (rollouts) & 2,585 & 0.032 & 65\% & 32.7\% & 82.3\% & 14.3 \\
\addlinespace[2pt]
\midrule
\multicolumn{7}{l}{closed-loop axis: measured on the four RoboCasa tasks
 (Table~\ref{tab:threebythree}); \pending{not yet measured on the other 20}} \\
\bottomrule
\end{tabular}
\caption{Every dataset we labeled, 24 in total. $\delta$ is the automatic
per-task threshold, unstable is the share of stamps with $\lambda>\delta$, and
deadband is $\lvert\lambda\rvert\le\delta$ where the label is ambiguous by
construction. We omit the stable share because the threshold rule fixes it near
2.5\% on every dataset, so it carries no information. The last two columns
measure heterogeneity. Within-ep var is the share of the total variance of
$\lambda$ that survives removing the between-episode means, so a value near
100\% means stability varies inside episodes rather than between them; it is 82
to 99.6\% on every dataset, which is the central evidence that chunk length
should not be a per-task constant. Crossings/100 counts regime transitions per
100 labeled stamps, normalised so longer episodes do not inflate it: the kitchen
tasks sit at 2 to 5 and alternate rarely, while long-horizon and high-precision
tasks sit at 10 to 25.}
\label{tab:alllabels}
\end{table*}

\begin{figure}[t]
\centering
\scriptsize
\begin{tabular}{@{}ll@{}}
\multicolumn{2}{@{}l}{\textit{RoboCasa PickPlaceToCabinet}, $\delta{=}.065$} \\
\quad demo 52, 6 crossings & \texttt{..................................................UUU...} \\
\quad demo 152, 26 crossings & \texttt{..........U................................U............} \\
\addlinespace[3pt]
\multicolumn{2}{@{}l}{\textit{RoboCasa TurnOnSinkFaucet}, $\delta{=}.058$} \\
\quad demo 131, 2 crossings & \texttt{............................................U...........} \\
\quad demo 68, 26 crossings & \texttt{..U...U.U.UU.UUUU...........U........................UU.} \\
\addlinespace[3pt]
\multicolumn{2}{@{}l}{\textit{LIBERO K8, both moka pots}, $\delta{=}.014$} \\
\quad demo 15, 43 crossings & \texttt{......UU...UUUUU......UUUUUUU.UUU..UUUUUUU....UUUUU.UUUU} \\
\quad demo 33, 73 crossings & \texttt{..UU....UU.UUUUUU.....U..U..U.U.UU.U.U.UUUU....U......UU} \\
\addlinespace[3pt]
\multicolumn{2}{@{}l}{\textit{robomimic tool\_hang}, $\delta{=}.035$} \\
\quad demo 42, 43 crossings & \texttt{.....UU..............U.UU..UUU........UU..U.UUU.UUU.UUUU} \\
\quad demo 156, 79 crossings & \texttt{.........UUUUU..U.U...U.UU..U.UU......UUUUU...UUUU.U.U.U} \\
\end{tabular}
\caption{Regime timelines for individual demonstrations, each row one episode
resampled to 56 columns. A dot is a stable or deadband stamp and the letter U is an unstable one. The kitchen tasks put one or two short unstable islands in a
long stable episode. The long-horizon LIBERO task and the high-precision
tool-hang task alternate throughout, which is the regime where a shorter
constant chunk, not switching, is the right answer.}
\label{fig:timelines}
\end{figure}

The unstable segments are short and localized. Measured on the states the
policy itself visits, a contiguous unstable stretch lasts a median of 1 to 5
control steps and occupies 1 to 3 percent of an episode on three of the four
tasks. OpenDrawer is the exception, where the signal stays up for long
stretches covering about 18 percent of an episode, and we return to that task
as the case where switching degenerates.

\subsection{The structure is a property of the dynamics, not of the probe}

A reader may ask whether the regimes are an artifact of the perturbation size
we inject. They are not. Relabeling all four tasks at $\sigma_u$ of 0.025 and
0.100, a fourfold range around the 0.05 used throughout, leaves the per-stamp
ordering of $\lambda$ almost unchanged, with Spearman correlation 0.95 to 0.98
between adjacent settings and 0.82 to 0.90 between the extremes, and the
unstable flag agrees on 97 to 99 percent of stamps.

\subsection{Fixed chunk length follows the regime profile}

Table~\ref{tab:fixed} gives the fixed-$k$ sweep. The best constant is 16, 16, 8
and 16 across the four tasks, so no single value is right even inside one
benchmark. TurnOnSinkFaucet is the extreme case, collapsing from .42 to .04 as
$k$ shrinks, which is the signature of a task that cannot tolerate constant
replanning. A second policy class shows the same task dependence with a
different profile. ACT policies trained on the same demonstrations score
.58/.64/.52/.32 on OpenCabinet, .30/.32/.04/.10 on OpenDrawer,
.12/.12/.04/.10 on PickPlaceCounterToCabinet and .50/.40/.34/.36 on
TurnOnSinkFaucet for $k$ of 16, 8, 4 and 1, so ACT tolerates single-step
execution on the faucet task where the generalist policy does not.

\begin{table}[t]
\centering
\small
\begin{tabular}{lcccc}
\toprule
Task & $k{=}16$ & $k{=}8$ & $k{=}4$ & $k{=}1$ \\
\midrule
OpenCabinet & \textbf{.52} & .52 & .44 & .48 \\
OpenDrawer & \textbf{.64} & .56 & .64 & .40 \\
PickPlaceCounterToCabinet & .58 & \textbf{.78} & .70 & .52 \\
TurnOnSinkFaucet & \textbf{.42} & .36 & .10 & .04 \\
\bottomrule
\end{tabular}
\caption{RoboCasa with GR00T N1.5, success over 50 episodes per cell. Bold
marks the best constant chunk length per task.}
\label{tab:fixed}
\end{table}

\section{Claim 2: The Regime Can Be Labeled and Predicted}
\label{sec:claim2}

\subsection{Three ways to answer the same question}

Every controller in this paper answers one yes or no question at each replan:
is the chunk the policy just proposed open-loop unstable? If the answer is no,
the controller executes all 16 actions without looking. If the answer is yes,
it switches to single-step execution and replans after every action until the
signal clears. The controllers differ only in how they reach that answer.

The \textbf{simulator oracle} measures the answer instead of predicting it. At
each replan it saves the simulator state, replays the proposed chunk once
unperturbed and eight more times with the first action perturbed, fits
$\lambda$, restores the state exactly, and compares $\lambda$ to $\delta$. It
is trained on nothing. It cannot run on a real robot because it rolls the
simulator forward and then rewinds it, so its role is to bound what a perfect
answer would be worth.

The \textbf{h1 predictor} is a small attentive head on frozen V-JEPA features
that predicts $\lambda$ from one clip of 16 consecutive camera frames, about
0.8 seconds, plus the robot state over the same window. It is trained by
supervised regression and classification against the labels of
Table~\ref{tab:labels}, sees no simulator at run time, and costs one forward
pass per replan.

The \textbf{h4 predictor} has identical architecture and identical training
targets but reads four clips spaced 16 frames apart, about 3.2 seconds of
history, with a learned per-clip position embedding.

\subsection{Perception recovers the regime}

Table~\ref{tab:claim2} gives detection quality on held-out demonstrations. Both
heads rank unstable segments well above chance on all four tasks. The result
holds on the other platforms as well: 0.94 to 0.96 AUROC on confident stamps
for MimicGen coffee preparation, 0.83 to 0.94 for nut assembly, 0.88 to 0.91
across LIBERO tasks, and 0.87 on a pooled four-task robomimic set, with
cross-task transfer from square to tool hang at 0.72 to 0.77.

\begin{table}[t]
\centering
\small
\setlength{\tabcolsep}{4pt}
\begin{tabular}{llcccc}
\toprule
& & \multicolumn{2}{c}{AUROC, confident stamps} & \multicolumn{2}{c}{recall at matched budget} \\
\cmidrule(lr){3-4} \cmidrule(lr){5-6}
Platform & Task & h1 & h4 & h1 & h4 \\
\midrule
RoboCasa GR00T & OpenCabinet & \textbf{0.885}\,$\pm$0.021 & 0.858\,$\pm$0.014 & 0.33 & 0.17 \\
RoboCasa GR00T & OpenDrawer & \textbf{0.863}\,$\pm$0.023 & 0.838\,$\pm$0.010 & 0.18 & 0.10 \\
RoboCasa GR00T & PickPlaceToCabinet & \textbf{0.895}\,$\pm$0.016 & 0.887\,$\pm$0.006 & 0.25 & 0.14 \\
RoboCasa GR00T & TurnOnSinkFaucet & \textbf{0.902}\,$\pm$0.022 & 0.895\,$\pm$0.012 & 0.20 & 0.15 \\
MimicGen DP & coffee preparation & 0.939\,$\pm$0.006 & \textbf{0.953}\,$\pm$0.002 & -- & -- \\
MimicGen DP & nut assembly & \textbf{0.885}\,$\pm$0.029 & 0.859\,$\pm$0.022 & -- & -- \\
\midrule
\multicolumn{6}{l}{\textit{other evidence, single-clip architecture}} \\
LIBERO-10 DP & 10 tasks & \multicolumn{2}{c}{.843 to .913} & \multicolumn{2}{c}{--} \\
robomimic DP & 4-task pool & \multicolumn{2}{c}{.868} & \multicolumn{2}{c}{--} \\
robomimic DP & square $\to$ tool\_hang & \multicolumn{2}{c}{.719 to .768} & \multicolumn{2}{c}{--} \\
robomimic DP & DINOv2, single frame & \multicolumn{2}{c}{.859} & \multicolumn{2}{c}{--} \\
any & privileged contact flag & \multicolumn{2}{c}{.602 to .802} & \multicolumn{2}{c}{--} \\
\midrule
\multicolumn{6}{l}{h1 and h4 on the remaining 16 labeled datasets:
 \pending{in the training pipeline}} \\
\multicolumn{6}{l}{pooled head over all four RoboCasa tasks:
 \pending{in the training pipeline}} \\
\multicolumn{6}{l}{head trained on both stability axes jointly:
 \pending{not started, waiting on closed-loop labels}} \\
\bottomrule
\end{tabular}
\caption{Detection quality of the stability predictors. AUROC is measured on
held-out demonstrations, mean and standard deviation over three seeds, and
restricted to confident stamps whose label is not inside the deadband. Recall
at matched budget is the fraction of truly unstable stamps recovered when the
firing rate is fixed to the label base rate, where random firing would recover
2.5 to 5.0 percent, so the heads run 3 to 13 times better than chance. Bold
marks the better of h1 and h4 per row: extra temporal history helps on one
dataset out of six and hurts on the other five. The privileged contact flag,
which reports whether the gripper touches a movable object, is the baseline a
learned head must beat to show it is doing more than contact detection.}
\label{tab:claim2}
\end{table}

The heads are not merely detecting contact. A privileged flag that reports
whether the gripper touches a movable object reaches only 0.60 to 0.80 AUROC on
the same stamps, so the learned heads discriminate which contact moments are
sensitive rather than that contact is occurring.

Whether temporal history helps depends on the task, and the dependence is
systematic. On the four RoboCasa tasks the four-clip head is worse than the
single-clip head everywhere, which suggested at first that history is simply
unhelpful. Extending the comparison to 14 datasets reverses that reading: the
single-clip head wins on 6 and the four-clip head on 8, with mean AUROC 0.848
against 0.858. What separates them is how fast the task alternates between
regimes. Ranking datasets by crossings per 100 stamps from
Table~\ref{tab:alllabels}, the advantage of history correlates with the
alternation rate at Spearman $+0.64$: on the seven slowest-alternating
datasets the four-clip head loses by 0.013 on average and wins 2 of 7, while
on the seven fastest it gains 0.034 and wins 6 of 7.

We read this as follows. A single frame suffices when instability arrives as
isolated islands, because the regime is then a property of the current
configuration, whether the gripper is engaged and how the geometry sits. A
single-frame DINOv2 head reaching 0.86 AUROC on confident stamps and
proprioception-only heads matching vision-only heads on MimicGen both support
that. When regimes alternate several times per hundred timesteps the current
configuration is ambiguous and recent motion carries the extra information.
RoboCasa, our main platform, sits firmly in the first case at 2 to 5 crossings
per 100 stamps, which is why the four-clip head loses there.
\pending{8 of 22 datasets still training; the correlation is over the 14 with
both heads complete}

\subsection{Existing replan signals do not track the measured one}
\label{sec:proxies}

Because our labels exist independently of any controller, they can be used to
audit the signals other methods switch on. We instrument shadow rollouts to
record, at the same state and the same replan, our $\lambda$ alongside the
sample spread of repeated chunk draws, the agreement between the queued chunk
and a freshly sampled one, and the speed profile of the predicted chunk, which
are the signals used by the confidence, consistency and kinematic families
respectively. Table~\ref{tab:proxies} reports the result over 1{,}370 to
2{,}339 replans per task.

\begin{table}[t]
\centering
\small
\begin{tabular}{lcccc}
\toprule
Signal family & Spearman with $\lambda$ & AUROC for $\lambda > \delta$ \\
\midrule
sample spread (confidence) & $-.07$ to $+.10$ & .61 to .76 \\
chunk agreement (consistency) & $-.17$ to $.00$ & .40 to .62 \\
predicted speed valley (kinematic) & $+.15$ to $+.21$, inverted & .31 to .69 \\
\bottomrule
\end{tabular}
\caption{Competing replan signals measured against $\lambda$ at the same
states, ranges across the four RoboCasa tasks. The kinematic row is inverted:
commanded speed correlates positively with $\lambda$, so firing at speed
valleys fires at the more stable moments.}
\label{tab:proxies}
\end{table}

None of the three tracks $\lambda$ closely. The confidence family carries a
weak tail signal but almost no graded correlation. The consistency family is
near chance. The kinematic family is the striking case: commanded speed
correlates positively with $\lambda$, so a rule that fires at speed valleys
fires preferentially at the more stable moments, the opposite of what the
stability measurement recommends. This does not say those methods do not work,
since a trigger can help for reasons unrelated to stability. It says they are
measuring something else, which is why our labels are not redundant with them.

\paragraph{Running them as controllers.}
Correlation says what a signal measures, not what it is worth. We therefore
reimplemented each signal as a controller in our own executor, switching
between $k{=}16$ and $k{=}1$ on the same policy, tasks and episode budget as
our own runs, with each threshold set so the signal fires at the oracle's rate
for that task. DEHP is excluded because its trigger is a horizon head trained
with reinforcement learning rather than a computable statistic.
Table~\ref{tab:baselines} gives the result.

\begin{table}[t]
\centering
\small
\setlength{\tabcolsep}{4pt}
\begin{tabular}{lcccc}
\toprule
Controller & signal family & mean & achieved firing & target \\
\midrule
simulator oracle & measured $\lambda$ & \textbf{.637} & .06 to .15 & matched \\
ours, h4 & learned $\lambda$ & .620 & .14 to .28 & matched \\
SGAC chunk cosine & consistency & .615 & .02 to .08 & \textbf{under} \\
DVAC sample variance & confidence & .585 & .06 to .15 & matched \\
BID chunk L2 & consistency & .585 & .00 to .09 & \textbf{under} \\
ours, h1 & learned $\lambda$ & .570 & .15 to .38 & over \\
PACE speed valley & kinematic & .555 & .07 to .14 & matched \\
\midrule
vanilla, fixed $k{=}16$ & none & .540 & 0 & \\
\bottomrule
\end{tabular}
\caption{Replan signals run as controllers on GR00T and RoboCasa, mean success
over the four tasks at 50 episodes each. Every controller uses the identical
executor and differs only in the statistic it thresholds. Achieved firing is
the range of realised firing rates across the four tasks. Thresholds were set
from a shadow pass in which the signal is scored but never acted on; once a
signal acts it changes the states the policy visits, so realised rates drift
from their targets. The two consistency signals drifted down, and one cell,
BID on OpenCabinet, fired on zero replans and is therefore a fixed-16 run in
disguise. Their means are consequently obtained with roughly half our
intervention budget, so those two rows are not yet like-for-like and we report
the realised rates rather than hide the discrepancy.
\pending{recalibrate the consistency signals against their live firing rates
and rerun}}
\label{tab:baselines}
\end{table}

Read with that caveat, the ranking is unremarkable: every learned or heuristic
controller lands within about six points of the others, and only the simulator
oracle separates from the pack. Nothing here supports a claim that our signal
makes a better trigger than the published alternatives at equal budget; what
it does support is that none of them, including ours, converts into a large
success gain on these tasks.

\subsection{Both stability axes measured at the same states}
\label{sec:threebythree}

Open-loop stability is one axis. The regime that would justify switching is
open-loop unstable together with closed-loop stable, and deciding membership
needs the closed-loop measurement of Section~\ref{sec:claim2}. We ran it
online on all four RoboCasa tasks, at every replan of a fixed-16 rollout, so
both axes are measured at the same states on the distribution the policy
actually visits. Table~\ref{tab:threebythree} is the result over 2{,}925
paired probes.

\begin{table}[t]
\centering
\small
\begin{tabular}{lrrrr}
\toprule
& \multicolumn{3}{c}{closed-loop} & \\
\cmidrule(lr){2-4}
open-loop & stable & deadband & unstable & row \\
\midrule
stable & 1.7\% & 8.7\% & 0.2\% & 10.6\% \\
deadband & 7.8\% & 47.4\% & 2.1\% & 57.3\% \\
unstable & \textbf{7.0\%} & 24.1\% & 1.0\% & 32.1\% \\
\midrule
column & 16.4\% & 80.2\% & 3.3\% & \\
\bottomrule
\end{tabular}
\caption{Both stability axes at the same states, pooled over the four
RoboCasa tasks, 2{,}925 paired probes from 95 episodes. The closed-loop value
is the growth rate of the perturbation minus the rate at which two
\emph{unperturbed} policy rollouts separate, so it measures what the
perturbation adds beyond the policy's own sampling. Both axes use the same
per-task threshold, so they ask the same question on the same scale. The bold
cell is the regime that motivates switching.}
\label{tab:threebythree}
\end{table}

Three readings matter. The switching regime exists but is a minority: 7.0
percent of all stamps, which is 22 percent of the open-loop-unstable stamps.
Its size tracks each task's own instability, running 11, 16, 23 and 31 percent
of the unstable stamps on OpenCabinet, OpenDrawer,
PickPlaceCounterToCabinet and TurnOnSinkFaucet. Closed-loop instability is
rare everywhere, 1.0 percent pooled and at most 5.6 percent on any task, so on
these tasks the four-regime picture is closer to a two-regime one.

The control branch was necessary rather than decorative. On every task, at
about 66 percent of stamps the perturbed closed-loop branch grows
\emph{more slowly} than two unperturbed rollouts do, meaning the policy's own
sampling separates trajectories faster than our injected perturbation. Without
subtracting that floor the closed-loop axis measures diffusion sampling rather
than stability, which is the defect in an earlier closed-loop labeling we
discarded.

This also sets a ceiling on what any open-loop detector can do. Our predictors
are trained on the open-loop axis alone, so even perfect open-loop detection
spends about 78 percent of its interventions at moments where closing the loop
cannot help. That is a quantitative reason to expect the intervention results
of Section~\ref{sec:claim3} to be weak, without the stability signal itself
being wrong. Training a head on the conjunction of both axes is the direct
consequence and we leave it to future work.

\section{Claim 3: Setting Chunk Length from the Predicted Regime}
\label{sec:claim3}

This is the claim for which our evidence is incomplete, and we report it as
such.

\pending{SAMPLE SIZE, to fix before submission. Every cell in this section is
50 episodes at a single seed, which gives a standard error near 7 points on a
success rate and near 10 points on the difference between two cells. Published
adaptive-chunking work evaluates far more heavily: PACE reports 900 episodes
per task-method pair over 50 tasks, DEHP reports 1000 environments over 3
seeds, SGAC averages 100 episodes over 3 seeds. Our present numbers therefore
cannot resolve effects of the size those papers report, and the final version
must raise every cell in Tables~\ref{tab:main} and~\ref{tab:fork} to at least
3 seeds of 50 episodes, and preferably widen the task set, before any of these
comparisons is stated as a result.}

\subsection{Getting the threshold right is most of the problem}

A learned head needs a firing threshold, and the obvious choices both fail.
Thresholding the predicted $\lambda$ at the label-space $\delta$ never fires at
all: a regression head compresses its outputs toward the mean, so on
TurnOnSinkFaucet the predicted $\lambda$ peaks at .064 against a $\delta$ of
.058 and only 0.12 percent of held-out predictions exceed it. Every switching
run built this way executes fixed 16 with extra computation. Calibrating the
threshold on held-out demonstrations instead fixes the scale but not the
distribution, and produces online firing rates from 2 to 29 percent against a 4
percent target, because the policy visits states unlike the demonstrations. We
therefore calibrate on deployment: a shadow run scores every replan without
acting, and the threshold is set at the quantile of that online distribution
which reproduces the oracle's firing rate for the task. Each predictor then
receives the same intervention budget as the oracle. Neither failure is visible
in AUROC, which is unchanged across all three choices.

\subsection{Task success}

Table~\ref{tab:main} is the main result of this section.

\begin{table*}[t]
\centering
\footnotesize
\setlength{\tabcolsep}{4.5pt}
\begin{tabular}{lccccccccc}
\toprule
& \multicolumn{3}{c}{No stability signal} & \multicolumn{1}{c}{Measured} & \multicolumn{4}{c}{Learned, no simulator} \\
\cmidrule(lr){2-4} \cmidrule(lr){5-5} \cmidrule(lr){6-9}
Task & vanilla & always & best fixed & oracle & h1 & fired & h4 & fired \\
& $k{=}16$ & $k{=}1$ & ($k$) & $(16,1)$ & $(16,1)$ & & $(16,1)$ & \\
\midrule
OpenCabinet & .52 & .48 & .52 (16) & .60 & .54 & .16 & .60 & .28 \\
OpenDrawer & .64 & .40 & .64 (16) & .62 & .60 & .34 & .62 & .26 \\
PickPlaceCounterToCabinet & .58 & .52 & \textbf{.78} (8) & .77 & .60 & .38 & .66 & .14 \\
TurnOnSinkFaucet & .42 & .04 & .42 (16) & .56 & .54 & .15 & \textbf{.60} & .14 \\
\midrule
mean & .54 & .36 & .59 & .64 & .57 & & .62 & \\
\bottomrule
\end{tabular}
\caption{RoboCasa with GR00T N1.5, success over 50 episodes per cell. The pair
$(16,1)$ under a switching controller gives its two chunk lengths: it executes
the long chunk $k_s{=}16$ when the segment is judged open-loop stable and drops
to the short chunk $k_u{=}1$ when it is judged unstable. The oracle column is
the mean over available seeds, two for PickPlaceCounterToCabinet and
TurnOnSinkFaucet and one otherwise. The fired columns give the fraction of
replans at which the predictor chose the short chunk. The always $k{=}1$ column
is the opposite extreme from the vanilla policy. The best fixed column names
the constant chunk length that achieves it in parentheses, from
Table~\ref{tab:fixed}; OpenCabinet ties at $k{=}8$ and OpenDrawer at $k{=}4$
and we report the longer chunk. Repeated runs of an identical configuration
differ by about .10 at this sample size.}
\label{tab:main}
\end{table*}

Three readings, and one caution.

The learned predictors improve on the vanilla policy on average, by 3 points
for h1 and 8 points for h4, and they never collapse. The h4 mean of .62 closes
most of the gap between the vanilla policy at .54 and the simulator oracle at
.64, which is the strongest single number in favour of the approach here. The always $k{=}1$ column
shows that this is not automatic: replanning every step for the whole episode
averages .36, eighteen points below the vanilla policy, and reaches .04 on
TurnOnSinkFaucet. The predictors enter that same single-step mode but only at
selected moments and stay above the vanilla policy everywhere, so a rare and
well placed switch is not the same intervention as a short fixed chunk.

The history predictor beats the single-clip predictor online on all four
tasks, reversing their offline ranking in Table~\ref{tab:claim2}. The firing
columns explain it. The two h1 cells that fire most, at .38 and .34, are the
ones that gain least, and across the eight learned cells the correlation
between firing rate and gain over the vanilla policy is $-0.65$. The online
advantage is calibration, not detection.

No controller reaches the best fixed chunk length on
PickPlaceCounterToCabinet, where $k{=}8$ alone scores .78 and even the oracle
reaches .77. On the three tasks whose best constant is already 16, the adaptive
controllers gain a few points. Against a .10 run-to-run spread only
TurnOnSinkFaucet with h4 is clearly outside noise.

The caution is that a success rate cannot separate the effect of an
intervention from the difference between two independent sets of random
episodes. The next experiment does separate them.

\subsection{A paired test of the intervention}

We fork each episode at the decision point. An episode runs under the oracle
until the first replan where $\lambda$ crosses $\delta$. At that point we
snapshot the simulator and run the remainder twice from that identical state,
once with the oracle acting and once with the oracle ignored so the policy
continues with fixed 16. The branches share their entire history and differ
only in the decision at the fired moment, which removes episode-to-episode
variance completely.

\begin{table}[t]
\centering
\small
\begin{tabular}{lccccc}
\toprule
Task & forked & rescued & broken & unchanged & $p$ \\
\midrule
PickPlaceCounterToCabinet & 40 & 4 & 4 & 32 & 1.00 \\
TurnOnSinkFaucet & 22 & 0 & 0 & 22 & 1.00 \\
\midrule
combined & 62 & 4 & 4 & 54 & 1.00 \\
\bottomrule
\end{tabular}
\caption{Paired fork test. Rescued counts episodes that succeed only when the
loop is closed at the fired moment, broken counts those that succeed only when
it is ignored, and $p$ is an exact two sided McNemar test. Success rates among
forked episodes are identical in both branches, .625 and .455.}
\label{tab:fork}
\end{table}

The result is a clean null (Table~\ref{tab:fork}). Closing the loop at a
measured unstable moment rescues 6.5 percent of forked episodes and breaks an
equal 6.5 percent. A paired comparison of time to success among episodes that
succeed in both branches is also null and slightly favours ignoring the signal.
The branches do diverge strongly, by up to a factor of five in steps to
completion, so the null is not an artifact of the two arms staying identical.

This also explains the pattern in Table~\ref{tab:main}. The largest oracle
number, .77 on PickPlaceCounterToCabinet, sits below a fixed chunk of 8 at .78,
and a second seed of the faucet oracle gives .50 where the first gave .62. What
looked like a causal gain in single unpaired cells does not survive pairing.

\subsection{The same conclusion on a second policy class}

We ran the equivalent sweep with diffusion policies on LIBERO, MimicGen and
robomimic: 42 predictor-driven cells on tasks that carry signal. Only 9 beat
their fixed-16 baseline, and none beats the best constant chunk for its task.
Where a shorter constant helps, using it everywhere is better than switching:
nut assembly reaches .52 at $k{=}1$ against a best predictor cell of .46, and
three-piece assembly reaches .78 at $k{=}1$ against .62. The damage is
graded by how often the predictor fires, with a correlation of $+0.58$ between
mean executed chunk length and change against baseline. Cells that fire most
lose 23 points on average, cells that fire least are neutral. Those runs used
demonstration-calibrated thresholds, so part of that damage is the calibration
failure described above, but even the lightly firing cells average exactly
zero benefit.

\subsection{A composed task with a built-in regime change}

Every task above is a single stage. If the hypothesis is right, switching
should pay most where an episode is guaranteed to cross regimes. We build such
an episode by forcing the cabinet shut at reset inside a
PickPlaceCounterToCabinet environment, so one 1800-step episode contains an
opening stage and then a pick-and-place stage, selected by the language
instruction. Stage 1 ends on the environment's own door-open predicate; full
success needs the environment's pick-and-place predicate as well.

\begin{table}[t]
\centering
\small
\begin{tabular}{lcc}
\toprule
Controller & full task & stage 1 \\
\midrule
fixed $k{=}16$ & .38 & .66 \\
fixed $k{=}8$ & .44 & .68 \\
fixed $k{=}4$ & .22 & .48 \\
fixed $k{=}1$ & .10 & .62 \\
\midrule
h1 $(16,1)$ & .48 & .66 \\
h4 $(16,1)$ & .48 & \textbf{.78} \\
simulator oracle $(16,1)$ & \textbf{.54} & .79 \\
\bottomrule
\end{tabular}
\caption{Composed OpenCabinet then PickPlaceCounterToCabinet in one 1800-step
episode, 50 episodes per cell. This is the only setting in the paper where
switching beats the best constant chunk. The oracle cell is
\pending{28 of 50 episodes at the time of writing}. The learned cells select a
stage-matched head using the environment's stage predicate, which is
privileged information the other cells do not receive;
\pending{a pooled head over all four tasks, which needs no stage signal, is
training and will replace these cells}.}
\label{tab:composed}
\end{table}

Both learned predictors reach .48 against .44 for the best constant chunk, and
the oracle is higher still at .54. The stage-1 column locates the effect: the
regime-aware controllers get through the opening stage more often than any
constant chunk, .78 and .79 against .68. This is the result the hypothesis
predicts and the only place we observe it, which we read as support for the
mechanism and against the generality of the intervention rather than the
reverse.

Two caveats keep it from being a headline. The learned cells receive the stage
boundary from the environment, and a pooled head trained over all four tasks,
which needs no such signal, costs real accuracy: per-task AUROC falls to .69
to .80 from the .863 to .902 of single-task heads. Removing the privilege is
not free, and the honest version of this table is still running.

\subsection{Seed variance is larger than the effects}

Every cell above is a single 50-episode run. Repeating four configurations at
two further seeds gives the spread in Table~\ref{tab:seeds}, and it is
sobering.

\begin{table}[t]
\centering
\small
\begin{tabular}{llccc}
\toprule
Task & controller & seed 0 & seed 1 & spread \\
\midrule
PickPlace & vanilla & .58 & .62 & 4 \\
PickPlace & h1 & .60 & .70 & 10 \\
PickPlace & oracle & .80 / .74 & .60 & \textbf{20} \\
TurnOnSinkFaucet & vanilla & .42 & .54 & \textbf{12} \\
TurnOnSinkFaucet & h1 & .54 & .66 & 12 \\
TurnOnSinkFaucet & oracle & .62 / .50 & .48 & 14 \\
\bottomrule
\end{tabular}
\caption{Success across seeds for identical configurations, 50 episodes each,
spread in percentage points. The vanilla baseline alone moves 12 points on
TurnOnSinkFaucet and the PickPlace oracle moves 20, which is larger than any
effect reported in this section.
\pending{full 3-seed replication in progress, 8 of 32 cells}}
\label{tab:seeds}
\end{table}

The .77 pooled oracle figure on PickPlaceCounterToCabinet does not reproduce
at a third seed, and the vanilla baseline on TurnOnSinkFaucet moves further
between seeds than the gap we attribute to switching. We therefore treat every
single-seed number in this section as provisional, and none of the
differences in Table~\ref{tab:main} as established.

\subsection{Where this leaves Claim 3}

Regime-conditioned chunk length is, on current evidence, worth a few points
over a fixed 16-step chunk and not worth anything over a properly swept
constant. The intervention does not survive a paired causal test at the
operating point we tested, which is a drop to single-step execution. Three
things remain untested and are the natural next steps: milder switches such as
$k{=}4$ at the fired moments, forking on tasks other than the two reported
here, and the closed-loop axis, which is the axis that would tell us in advance
whether closing the loop at an unstable moment can help at all.
